\documentclass[10pt,a4paper]{article}

% ---------- PACKAGES ----------
\usepackage[margin=1.5cm]{geometry}
\usepackage{amsmath, amssymb}
\usepackage{physics}
\usepackage{enumitem}
\usepackage{multicol}
\usepackage{titlesec}
\usepackage{mathrsfs}
\usepackage{bm}
\usepackage{tikz}
\usepackage{mecaso}
\usepackage{xparse}

\usetikzlibrary{calc}

% ---------- MACROS PERSONNALISEES ----------

% Torseur symbolique avec point d'application facultatif
\NewDocumentCommand{\torseur}{o}{%
    \IfValueTF{#1}
        {\left\{\mathscr{T}\right\}_{#1}}
        {\left\{\mathscr{T}\right\}}%
}

% ---------- MISE EN FORME ----------
\setlist[itemize]{noitemsep, topsep=0pt}
\renewcommand{\vec}{\overrightarrow}

\author{Raphaël Jamann}
\title{Formulaire Mécanique Lagrangienne \\ 3GM}
\date{}

% ======================================================
% DOCUMENT
% ======================================================

\begin{document}

    \maketitle

    \begin{multicols}{2}

    % ======================================================
    \section{Actions mécaniques}

    \subsection*{Moment d'une force}

    \[
    \vec M_O(F) = \vec{OA} \wedge \vec F
    \]

    Avec A le point d'application de la force.

    \subsection*{Transport du moment (BABAR)}

    \[
    \vec M(B) = \vec M(A) + \vec{BA} \wedge \vec R
    \]


    \subsection*{Transport des Vitesses (formule de Varignon)}
        Pour appliquer la formule de Varignon, il faut que le point $A$ et $VB$ appartiennent au même solide $S$.
        $$ \VARK{A}{B}{S} $$

    % ======================================================


    % ======================================================
    \section{Dérivée d'un vecteur tournant}

    Cette formule est notamment très utile pour calculer la vitesse d'un point $P$ d'un solide $S_i$ par rapport au repère $R_0$.

    $$ \Din[0]{\Vct{OP}}=\Din[i]{\Vct{OP}}+\Rotation[0]{i}\times\Vct{OP} $$

    % ======================================================
    \section{Torseurs}

        \subsection*{Cinématique}

        Le mouvement d'un solide $S$ par rapport à $R$ est décrit par le torseur :

        \[ \Ktor[R]{A}{S} \]

        \subsection*{torseur d'actions mécaniques}

        $$ \Stor[0]{A}{S} $$

        $$ \LFtor[S/0]{F}{X}{Y}{Z}{L}{M}{N}{A} $$

        % ======================================================
        \subsection*{Cinétique}

        \[
        \Ctor[0]{A}{S}
        \]

        % ======================================================
        \subsection*{Dynamique}

        \[
        \Dtor[0]{A}{S}
        \]

    \end{multicols}

    \clearpage

    \section{Energie Cinétique}

        \subsection*{du point matériel P}
            $$ T_{P/0}  = \KineticT[0]{A}{S} $$
            Lorsqu'on calcule par exemple $\vec v^2 = (\dot{\rho} \vec{u_0} + \rho \dot{\theta} \vec{v_1})^2$, on utilise l'identité remarquable suivante :
            $$ \vec v^2 = \dot{\rho}^2 + (\rho \dot{\theta})^2 + 2\dot{\rho} \vec{u_0} \cdot \vec{v_1} \rho \dot{\theta} $$
            Ou alors plus généralement si $\vec v = a \vec u + b \vec v + c \vec w$ alors
            $$ \vec v^2 = a^2 + b^2 + c^2 + 2a \vec u \cdot \vec v b + 2a \vec u \cdot \vec w c + 2b \vec v \cdot \vec w c $$
        \subsection*{du solide}
            $$ T_{1/0} = \KEtor[0]{1} $$
            Si le point $A \equiv G$ centre de masse du solide :
            $$ T_{1/0} = \KineticG[0]{S} $$
            Si le point $A$ est fixe dans $R_0$ alors
            $$ T_{1/0} = \KineticR[0]{A}{S} $$

        \subsection*{d'un ensemble de solides}
            Si $\Sigma = \{S_1, S_2, \ldots, S_n\}$ alors
            $$ T_\Sigma = T_{1/0} + T_{2/0} + \cdots + T_{n/0} $$

    \section{Energie Potentielle}
        \subsection*{Pesanteur}

        $$ U_{\text{grav}} = -m_i \vec r_{G_i} \cdot \vec g + \text{cste} $$

        \subsection*{Ressort}
        Traction / Compression :
        $$ U_{\text{ressort}} = \frac{k}{2} (L- L_0)^2 + \text{cste} $$
        Torsion:
        $$ U_{\text{ressort}} = \frac{k}{2} (\theta - \theta_0)^2 + \text{cste} $$

    \section{Puissance}
        \subsection*{Puissance développée par un torseur des actions mécaniques extérieures au système sur un solide $S_i$}
            $$ \mathscr{P}_{\text{ext}\rightarrow S_i/0} = \Ptor[0]{S_i}{ext} $$

        \subsection*{Puissance développée dans une liaison entre deux solides $S_i$ et $S_j$ du système considérée}
            $$ \mathscr{P}_{S_i \leftrightarrow S_j} = \Ptor[i]{S_j}{S_i} $$
            $$ \mathscr{P}_{S_i \leftrightarrow S_j} = \mathscr{P}_{S_j \leftrightarrow S_i} $$

        \subsection*{Puisance développée par une force conservative}
            $$ \mathscr{P}_{\text{cons}\rightarrow S_i/0} = - \frac{dU}{dt} $$


    \clearpage

    \section{Equation de Lagrange}

        $$ \mathscr{L}(q_i): \qquad \frac{d}{dt} \left( \pdv{T_{\Sigma/0}}{\dot{q}_i} \right) - \pdv{T_{\Sigma/0}}{q_i} = Q_i = Q_i^{\text{c}} + Q_i^{\text{nc}} $$
        Or les forces généralisées conservatives peuvent se calculer avec $ Q_i^{\text{c}} = - \pdv{U_{\Sigma}}{q_i} $ donc
        $$ \mathscr{L}(q_i): \qquad \frac{d}{dt} \left( \pdv{T_{\Sigma/0}}{\dot{q}_i} \right) - \pdv{T_{\Sigma/0}}{q_i} + \pdv{U_{\Sigma}}{q_i} = Q_i^{\text{nc}} $$

        Pour trouver les forces généralisées non-consevatives (ou conservatives mais on a la formule pour aller plus vite), il faut calculer la puissance virtuelle puis identifier.
        Par exemple, si on a deux forces généralisées $Q_\alpha$ et $Q_\beta$ associées aux coordonnées généralisées $\alpha$ et $\beta$, on a:
        $$ \mathscr{P}^* = Q_{\alpha} \dot{\alpha}^* + Q_{\beta} \dot{\beta}^* $$

    % ======================================================
    \section{Principe Fondamental de la Dynamique}


        \subsection*{Principe Fondamental de la Statique}

        \[
        \sum \torseur[A] = \{0\}
        \]

        \subsection*{Théorème de la résultante dynamique}
        \textbf{TRD :}

        \[
        \sum \vec F_{\text{ext}}
        =
        m\vec A_G
        =
        \vec D
        \]

        \subsection*{Théorème du moment dynamique}
        \textbf{TMD :}

        \[
        \sum \vec{\mathcal M}
        =
        \vec\delta
        \]

\end{document}
